\section{Pixel-to-parameter mapping}
\label{sec:mapping}

Before any fractal iteration can begin, every pixel on screen must be
assigned a complex parameter~$c$.  At moderate zoom depths this mapping is
a trivial affine transform, but at the extreme magnifications targeted by
\FractalShark{}---$10^{10{,}000}$ and beyond---the pixel-to-parameter
conversion becomes a precision-critical operation.  Adjacent pixels may
differ by $10^{-\text{thousands}}$, yet the mapping must be accurate to
the last bit if the downstream perturbation pipeline is to produce correct
results.  The transform is therefore computed at full arbitrary precision
(\code{HighPrecision}) on the host, with only the final per-pixel offsets
projected down to the kernel's working type.

This section describes how CUDA threads index pixels, how the coordinate
system is oriented, and how the affine map is parameterized.

Each CUDA thread computes a pixel coordinate \((X,Y)\) and maps it to a complex
parameter \(c = x_0 + i y_0\). A typical kernel starts with:
\begin{verbatim}
int X = blockIdx.x * blockDim.x + threadIdx.x;
int Y = blockIdx.y * blockDim.y + threadIdx.y;
if (X >= width || Y >= height) return;
size_t idx = ConvertLocToIndex(X, height - Y - 1, width);
\end{verbatim}

\subsection{Thread-to-pixel mapping}
A 2D CUDA grid of 2D blocks covers the image. Each thread is responsible for one
pixel. The bounds check prevents out-of-range threads from writing.

\subsection{Y-axis convention}
The expression \code{height - Y - 1} flips \(Y\). This is common when the image
buffer uses a top-left origin but the complex-plane mapping assumes a
bottom-left origin (or vice versa).

\subsection{Affine map into the complex plane}
The parameter \(c\) is computed by an affine transform:

\begin{align}
x_0 &= cx + dx \cdot X, \\
y_0 &= cy + dy \cdot Y,
\end{align}

where \code{cx,cy} anchor the plane (e.g., the coordinate at pixel \((0,0)\)),
and \code{dx,dy} are per-pixel increments. Different kernels compute these
expressions in different numeric types; the intent is always the same: map each
pixel to its associated complex parameter \(c\).


\section{Mandelbrot recurrence in real arithmetic}
\label{sec:real-form}

Writing \(z = x + i y\) and \(c = x_0 + i y_0\), the iteration becomes:

\begin{align}
x_{n+1} &= x_n^2 - y_n^2 + x_0, \\
y_{n+1} &= 2 x_n y_n + y_0.
\end{align}

The escape test is:

\begin{equation}
x_n^2 + y_n^2 \ge 4 \quad\Rightarrow\quad \text{escaped.}
\end{equation}

Many kernels cache squares:
\[
zrsqr = x^2,\quad zisqr = y^2,
\]
so that \(x^2-y^2\) and \(x^2+y^2\) can be formed cheaply as \code{zrsqr - zisqr}
and \code{zrsqr + zisqr}. This is especially valuable when \(x\) and \(y\) are
represented by multi-component expansion types.


\section{Mandelbrot base kernels}
\label{sec:mandel-base-kernels}

The goal of these kernels in \FractalShark{} is to \emph{render the Mandelbrot set}
by computing an \emph{escape-time} iteration count (\cref{sec:mandelbrot-intro}) per pixel over a 2D image
grid efficiently. Each CUDA thread evaluates a single complex parameter \(c\)
corresponding to one pixel (\cref{sec:mapping}), iterates the Mandelbrot recurrence (\cref{sec:real-form}), and stores the
iteration count into an output buffer. Separate (or fused) stages can map
iteration counts to colors, apply palettes, and perform anti-aliasing
(\cref{subsec:coloring,subsec:antialiasing}).

Across these base kernels, the primary variation is the numeric representation
used for the orbit arithmetic: from IEEE-754 \code{float}/\code{double} up
through expansion types (float-float, double-double, quad-float, quad-double)
and HDR-normalized formats (\cref{sec:extended-precision-types}). The shared objective remains the same: compute the
escape-time for the Mandelbrot iteration as accurately and efficiently as needed
for a desired zoom depth.  A complete listing of all CUDA kernels is given in
\cref{sec:kernel-list}.

Each kernel variant renders the same Mandelbrot escape-time field; the only
difference is the numeric type used for mapping and orbit iteration. The
following sections describe how each type realizes the same recurrence and
escape test.  This set of kernels does not use linear approximation or
perturbation; they simply evaluate the Mandelbrot iteration directly in the
chosen numeric format.

Throughout, \code{IterType} is the integer type used to store the escape-time
iteration count (e.g., \code{uint32\_t} or \code{uint64\_t}).

\subsection{Common structure}

Every base kernel follows the same algorithmic template:
\begin{enumerate}
  \item Map pixel \((X,Y)\) to a complex parameter \(c = x_0 + i\,y_0\)
        via the affine mapping described in \cref{sec:mapping}.
  \item Initialize \(x = y = 0\) (or equivalently \(z_0 = 0\)).
  \item Iterate the real-arithmetic Mandelbrot recurrence (\cref{sec:real-form}):
        \[
          x \leftarrow x^2 - y^2 + x_0,\qquad y \leftarrow 2xy + y_0.
        \]
  \item After each iteration (or chunk of iterations), test \(x^2 + y^2 \ge 4\).
  \item Store the escape iteration count into the output buffer.
\end{enumerate}
The only variation across the seven kernel variants is the numeric type used for
all arithmetic and the associated precision--performance tradeoff.

\subsection{Kernel variants}

\Cref{tab:base-kernels} summarizes the seven base kernels. All double- and
quad-limb float/double expansion types are based on work by Andrew
Thall~\cite{andrew-thall-dblflt}, which builds on Dekker's
double-double technique~\cite{Dekker1971DoubleDouble}. The quad-component
types also draw on the QD library by Hida, Li, and
Bailey~\cite{Hida2001QD}. The implementations are modified in
\FractalShark{} to support both \code{float} and \code{double} base types.

\begin{table}[!htbp]
\centering
\small
\caption{Mandelbrot base kernel variants.}
\label{tab:base-kernels}
\begin{tabularx}{\textwidth}{l r X}
\hline
\textbf{Kernel} & \textbf{Bits} & \textbf{Type and notes} \\
\hline
\label{sec:mandel-1x-float}%
\code{mandel\_1x\_float}
  & 24
  & IEEE-754 \code{float}.
    Highest throughput; uses FMA intrinsics
    (\code{\_\_fmaf\_rd}) to reduce intermediate rounding. \\[4pt]
\label{sec:mandel-1x-double}%
\code{mandel\_1x\_double}
  & 53
  & IEEE-754 \code{double}.
    FP64 throughput is often much lower than FP32 on
    consumer GPUs. \\[4pt]
\label{sec:mandel-2x-float}%
\code{mandel\_2x\_float}
  & ${\sim}48$
  & \code{dblflt} (\cref{subsec:ep-dblflt}), a float-float expansion
    ($a_\mathrm{hi}+a_\mathrm{lo}$).
    Leading-component escape test for speed. \\[4pt]
\label{sec:mandel-2x-double}%
\code{mandel\_2x\_double}
  & ${\sim}106$
  & \code{dbldbl} (\cref{subsec:ep-dblflt}), a double-double expansion.
    Leading-component escape test for speed. \\[4pt]
\label{sec:mandel-4x-float}%
\code{mandel\_4x\_float}
  & ${\sim}96$
  & \code{GQF::gqf\_real} (\cref{subsec:ep-higher-prec}), a four-float
    expansion ($a_0+a_1+a_2+a_3$).
    Full-precision escape test. \\[4pt]
\label{sec:mandel-4x-double}%
\code{mandel\_4x\_double}
  & ${\sim}212$
  & \code{GQD::gqd\_real} (\cref{subsec:ep-higher-prec}), a four-double
    expansion.
    Full-precision escape test. Deepest non-HDR zoom. \\[4pt]
\label{sec:mandel-hdr-float}%
\code{mandel\_hdr\_float}
  & ${\sim}48$
  & \code{HDRFloat<CudaDblflt<dblflt>>} (\cref{subsec:ep-cudadblflt}).
    HDR wrapper around float-float expansion; uses
    \code{HdrReduce()} normalization and a reduced comparator
    for escape test. Primarily a testbed. \\
\hline
\end{tabularx}
\end{table}

\paragraph{FMA intrinsics.}
The \code{mandel\_1x\_float} kernel illustrates how fused multiply-add can
implement the recurrence directly:
\begin{verbatim}
ytemp = __fmaf_rd(-y, y, x0);     // x0 - y^2
xtemp = __fmaf_rd(x, x, ytemp);   // x^2 - y^2 + x0
xtemp2 = 2.0f * x;
y = __fmaf_rd(xtemp2, y, y0);     // 2xy + y0
x = xtemp;
\end{verbatim}
The \code{\_rd} suffix selects round-downward; \code{\_rn} (round-to-nearest)
or plain \code{fmaf} may be substituted when strict rounding mode is not required.

\paragraph{Leading-component escape test.}
The two-limb expansion kernels (\code{mandel\_2x\_float},
\code{mandel\_2x\_double}) compare only the leading component
(\code{head}) against the bailout:
\[
  zrsqr.\mathrm{head} + zisqr.\mathrm{head} < 4.
\]
This is fast but can misclassify points extremely near the boundary; a fully
robust test would incorporate both components.

\paragraph{HDR normalization and reduced comparator.}
The \code{mandel\_hdr\_float} kernel calls \code{HdrReduce()} on intermediate
values so that norms and comparisons operate on normalized representations.
Escape is tested via:
\begin{verbatim}
while (zsq_sum.compareToBothPositiveReduced(Four) < 0)
\end{verbatim}
This keeps comparisons meaningful and stable in HDR arithmetic.


\section{Iteration chunking via \code{iteration\_precision}}
\label{sec:chunking}

A few of the base kernels just described (\cref{sec:mandel-base-kernels}), which exclude linear approximation or
perturbation, are templated on an integer \code{iteration\_precision}
(\(1,2,4,8,16\)) and unroll multiple Mandelbrot steps inside the loop:

\begin{itemize}
  \item Each loop iteration performs \code{iteration\_precision} updates.
  \item The counter \code{iter} increases by that amount.
  \item The maximum iteration \code{n\_iterations} is adjusted so \code{iter}
        does not exceed the requested limit.
\end{itemize}

This approach reduces loop overhead. The trade-off is that the escape predicate
is typically checked only once per chunk; therefore the reported escape
iteration can be larger than the true first-escape iteration by up to
\code{iteration\_precision - 1}. For strict escape iteration counts (e.g., for
continuous/smooth coloring based on the first bailout), use a chunk size of 1 or
insert bailout checks within the unrolled body.

This optimization was mainly explored for educational purposes; in practice, the
benefit is small compared to other optimizations such as perturbation (\cref{sec:perturbation-concept}) and linear
approximation (\cref{sec:approx-accel}).
